% =============================================================
% บทที่ 5 - สรุปผลการดำเนินงาน
% ปรับปรุงให้ตรงกับสถานะโค้ดจริง ณ ปัจจุบัน (ส่วนเว็บ: ตุลาคม 2569 / ส่วน AI, Trip, แชตบอต: ปรับตามโค้ดสาขา dev ณ 13 ตุลาคม 2569 แล้ว โมเดลที่ใช้งานจริงคือ gemini-3.8-flash ผ่าน Vertex AI ทุกฟังก์ชัน) แทนตัวเลขร่างเดิมจาก
% Research Me (1).docx ซึ่งบันทึกไว้ตั้งแต่ช่วงที่ระบบยังแยกส่วนกันอยู่ (ยังไม่เชื่อมต่อ AI จริง)
% หมายเหตุ: หัวข้อ 5.1 สรุปผลเทียบกับวัตถุประสงค์ในบทที่ 1 (ตาราง tab:objectives) แทนตาราง "สิ่งที่ทำสำเร็จแล้ว/สิ่งที่ยังคงค้าง"
% รายโมดูลเดิม ซึ่งเป็นรูปแบบรายงานความก้าวหน้า งานที่ค้างของแต่ละส่วนย้ายไปอยู่ในหัวข้อข้อจำกัดและข้อเสนอแนะ
% =============================================================

% -----------------------------------------------------------
\section{สรุปผลการดำเนินโครงงาน}
% -----------------------------------------------------------

โครงงานนี้พัฒนาเว็บแอปพลิเคชันไดโนจนระบบทั้งเว็บแอปพลิเคชัน, ระบบหลังบ้าน, และ AI Service เชื่อมต่อและทำงานร่วมกันได้จริงแบบครบวงจร และนำขึ้นใช้งานบนเซิร์ฟเวอร์จริงแล้ว ทั้งฟังก์ชันแชตบอตที่ตอบกลับแบบทยอยแสดงผล ฟังก์ชันสร้างและตรวจสอบคุณภาพแผนการเดินทางด้วยกลไก LLM-as-a-Judge อัตโนมัติ และฟังก์ชันเสริมด้าน AI อื่น ๆ เช่น การดึงข้อมูลกิจกรรมจากข้อความโพสต์ และการสร้างคำอธิบายสถานที่เฉพาะราย ส่วนของเว็บแอปพลิเคชันประกอบด้วยระบบบัญชีผู้ใช้ ระบบสแกน QR และแลกพอยท์ ระบบรายงานข้อมูล และระบบจัดการหลังบ้านที่เชื่อมต่อฐานข้อมูลจริง ผลการประเมินความพึงพอใจด้วยแบบสอบถามจากผู้ใช้ \surveyN{} คนได้ค่าเฉลี่ยรวม \surveyOverallMean{} จาก 5 คะแนน อยู่ในระดับ\surveyOverallLevel{} (ตารางที่ \ref{tab:survey_items}) เมื่อเทียบกับวัตถุประสงค์ทั้ง 3 ข้อในบทที่ 1 ได้ผลดังตารางที่ \ref{tab:objectives}

% ค่าเฉลี่ยรายด้านในตารางนี้พิมพ์ด้วยมือจากตารางที่ \ref{tab:survey_items} (ข้อมูลชุด 32 คน)
% ถ้ารันสคริปต์แบบสอบถามใหม่ต้องแก้ตัวเลขในตารางนี้ด้วย
\begin{table}[H]
\centering
\caption{สรุปผลการดำเนินงานเทียบกับวัตถุประสงค์}
\label{tab:objectives}
{\small
% ชิดซ้ายเพราะคอลัมน์แคบ การจัดเต็มบรรทัดทำให้ช่องไฟระหว่างคำยืดห่าง
\begin{tabular}{|>{\raggedright\arraybackslash}p{3.1cm}|>{\raggedright\arraybackslash}p{4.5cm}|>{\raggedright\arraybackslash}p{4.3cm}|>{\raggedright\arraybackslash}p{1.7cm}|}
\hline
\textbf{วัตถุประสงค์} & \textbf{ผลการดำเนินงาน} & \textbf{ผลการประเมิน} & \textbf{สถานะ} \\
\hline
1) พัฒนาเว็บแอปพลิเคชันรวมศูนย์ข้อมูลแหล่งท่องเที่ยว ร้านอาหาร และกิจกรรมของจังหวัดขอนแก่น & เปิดใช้งานจริงที่ dinokhonkaen.app มีหน้ารายละเอียดสถานที่และกิจกรรม การค้นหา บัญชีผู้ใช้ ระบบสะสมและแลกพอยท์ และ Admin Portal สำหรับดูแลข้อมูล & แบบสอบถาม \surveyN{} คน: ด้านการทำงานของระบบ 4.22 ด้านความง่ายในการใช้งาน 4.56 และด้านการออกแบบและประสิทธิภาพ 4.48 อยู่ในระดับมากที่สุดทั้งสามด้าน & บรรลุ \\ \hline
2) พัฒนาแชตบอตให้บริการข้อมูล ตอบข้อซักถาม และให้คำแนะนำด้านการท่องเที่ยว & แชตบอตแบบ RAG ที่ค้นข้อมูลสถานที่ ฐานความรู้ และกิจกรรมพร้อมกัน แนะนำสถานที่ตามพื้นที่ที่ผู้ใช้ระบุ ค้นเพิ่มจาก Google เมื่อข้อมูลในพื้นที่ไม่พอ ตอบกลับแบบทยอยแสดงผล และสนทนาต่อเนื่องหลายรอบได้ & ด้านแชตบอต AI 4.02 อยู่ในระดับมาก โดยข้อความตรงกับคำถามของคำตอบได้ 4.14 และข้อความรวดเร็วในการตอบได้ 3.87 & บรรลุ \\ \hline
3) พัฒนาระบบวางแผนการท่องเที่ยวที่จัดลำดับแผนให้สอดคล้องกับความต้องการของผู้ใช้ & LLM คัดเลือกสถานที่ตามความสนใจ ประเภทอาหาร งบประมาณ และระยะเวลา ตัวจัดตารางเชิงกำหนดแน่นอนจัดลำดับและเวลา รวมถึงวันสุดท้ายที่จบที่สนามบินหรือสถานีขนส่ง และ LLM-as-a-Judge ตรวจคุณภาพแผน ผู้ใช้บันทึกและจัดการแผนได้ (UC-07) & ทดลอง Prompt 96 รัน ผ่านตัวชี้วัด Environment, Hard และ Commonsense ร้อยละ 100 (ตารางที่ \ref{tab:prompt_eval}) และด้านการวางแผนเที่ยวด้วย AI 4.08 อยู่ในระดับมาก & บรรลุ \\ \hline
\end{tabular}
\par}
\end{table}

ทั้งนี้ ด้านที่ผู้ใช้ให้คะแนนต่ำกว่าด้านอื่นคือส่วนที่เรียกใช้ LLM ทั้งแชตบอตและการวางแผนเที่ยว โดยเฉพาะความเร็วในการตอบสนองและความเป็นไปได้ของตารางเวลาในแผน ข้อจำกัดของระบบและงานที่ยังไม่ได้ดำเนินการ เช่น การทดสอบยอมรับจากผู้ใช้แบบกำหนดภารกิจ (UAT) สรุปไว้ในหัวข้อถัดไปและหัวข้อข้อเสนอแนะ

% -----------------------------------------------------------
\section{ข้อจำกัดของระบบ}
% -----------------------------------------------------------

\subsection{ข้อจำกัดด้านการทำงาน}

ระบบมีข้อจำกัดด้านการทำงาน ดังนี้

\begin{enumerate}
\needspace{5\baselineskip}
\item มาตรการด้านความปลอดภัยยังไม่ครบถ้วน แม้ระบบจะตรวจสอบการยืนยันตัวตนและสิทธิ์ที่ฝั่งเซิร์ฟเวอร์แล้ว แต่ยังไม่มีนโยบาย Content Security Policy เต็มรูปแบบ (กำหนดเพียง frame-ancestors) เนื่องจากหน้าเว็บโหลดสคริปต์และรูปภาพจาก Google Maps และ Cloudinary การตรวจชนิดไฟล์อัปโหลดอาศัย mimetype ที่ฝั่งผู้ใช้ส่งมา และยังไม่มีการล็อกบัญชีหลังเข้าสู่ระบบผิดซ้ำนอกเหนือจากการจำกัดอัตราคำขอ
\item การดูแลโครงสร้างฐานข้อมูลต้องดำเนินการด้วยตนเอง ฐานข้อมูลที่ใช้งานจริงไม่มีประวัติ migration ของ Supabase หากสั่งนำ migration ทั้งหมดไปใช้ซ้ำจะประมวลผลไฟล์เดิมทั้งหมดอีกครั้ง จึงต้องนำไฟล์ใหม่ไปใช้ทีละไฟล์ และสคริปต์ปรับปรุงระบบ (deploy.sh) ไม่ได้ดำเนินการส่วนนี้ให้
\item ข้อมูลสถานที่จาก Google เป็นปัจจุบันตามรอบการซิงก์เท่านั้น ไม่ใช่แบบเรียลไทม์ แม้ระบบจะซิงก์อัตโนมัติได้วันละครั้ง แต่แต่ละสถานที่ที่ซิงก์ต้องเรียก Google Places API ซึ่งมีค่าใช้จ่ายตามจำนวนครั้ง จึงต้องจำกัดด้วยรอบการซิงก์และเพดานคำขอต่อเดือน (ค่าเริ่มต้น 900 คำขอ) ข้อมูลที่เปลี่ยนระหว่างรอบ เช่น เวลาเปิด-ปิดชั่วคราว จะยังไม่ปรากฏจนถึงรอบถัดไป นอกจากนี้ ความคืบหน้าของชุดงานนำเข้าตามพื้นที่เก็บในหน่วยความจำของ Backend หากเซิร์ฟเวอร์เริ่มทำงานใหม่ระหว่างนำเข้า สถานที่ที่นำเข้าแล้วยังอยู่ครบ แต่ผู้ดูแลระบบจะไม่เห็นผลสรุปของชุดงานนั้น
\item สถานที่ที่แชตบอตค้นเพิ่มจาก Google ยังไม่ผ่านการตรวจสอบของทีมงาน ระบบแจ้งผู้ใช้ไว้ในคำตอบ แต่ข้อมูลเช่นเวลาเปิด-ปิดและช่วงราคาอาจไม่ถูกต้อง และการค้นเพิ่มจำกัดไว้วันละ 100 ครั้ง เมื่อครบเพดานแล้ว คำถามที่ระบุพื้นที่ซึ่งฐานข้อมูลมีข้อมูลน้อยจะได้คำตอบเท่าที่มีในฐานข้อมูลเท่านั้น ส่วนคำถามที่อ้างถึงสถานที่สำคัญ (เช่น ``แถวเซ็นทรัล'') ไม่มีการค้นเพิ่ม เพราะที่อยู่จาก Google ไม่ระบุชื่อสถานที่สำคัญ
\item การแนะนำย่านที่พักสำหรับทริปที่ไม่ได้ระบุจุดเริ่มต้นคำนวณจากระยะทางแบบเส้นตรงระหว่างสถานที่ในแผน และแผนนั้นถูกจัดโดยเริ่มจากใจกลางเมือง หากผู้ใช้จองที่พักในย่านที่แนะนำแล้ว เวลาเดินทางจริงจะต่างจากแผน จนกว่าจะสร้างแผนใหม่โดยปักหมุดที่พัก
\item แผนการเดินทางปรับแก้ไม่ได้หลังสร้างเสร็จ ระบบบันทึกทุกแผนลงตาราง trips, trip\_days และ trip\_items โดยอัตโนมัติ และผู้ใช้ที่เข้าสู่ระบบเปิดดู เปลี่ยนชื่อ ทำสำเนา และลบแผนได้ (UC-07) แต่ฟังก์ชันปรับแผนด้วยการกดถูกใจ/ไม่ถูกใจและช่องข้อเสนอแนะที่เคยพัฒนาไว้ถูกถอดออกจากหน้าผู้ใช้แล้ว ผู้ใช้ที่ต้องการแผนต่างจากเดิมจึงต้องกรอกเงื่อนไขและสร้างแผนใหม่ทั้งชุด นอกจากนี้ แผนของผู้ที่ไม่ได้เข้าสู่ระบบถูกบันทึกเพื่อสถิติของผู้ดูแลระบบเท่านั้นและเปิดดูซ้ำไม่ได้ และบัญชีหนึ่งบันทึกแผนได้ไม่เกิน 50 แผน
\item การทดสอบส่วนเว็บยังไม่เป็นระบบอัตโนมัติเต็มรูปแบบ กรณีทดสอบเชิงหน้าที่ TC-01 ถึง TC-37 รันด้วยสคริปต์แยกต่างหากบนฐานข้อมูลชั่วคราวและยังไม่รวมอยู่ในที่เก็บโค้ด Backend มีชุดทดสอบอัตโนมัติ 127 กรณีแล้ว (ตารางที่ \ref{tab:auto_tests}) แต่ใช้ฐานข้อมูลจำลองและยังไม่ถูกรันใน GitHub Actions Frontend ไม่มีชุดทดสอบอัตโนมัติ และ GitHub Actions ตรวจฝั่ง Frontend เฉพาะการสร้างโปรแกรม นอกจากนี้โครงงานไม่ได้ทดสอบภาระงานเมื่อมีผู้ใช้พร้อมกันจำนวนมาก การจำกัดอัตราคำขอเก็บสถานะในหน่วยความจำของกระบวนการเดียว จึงเหมาะกับการติดตั้งบนเครื่องเดียวเท่านั้น
\item ระยะทางและเวลาเดินทางในแผนเป็นค่าประมาณ ระบบคำนวณระยะทางแบบเส้นตรง (Haversine) และใช้ความเร็วเฉลี่ยคงที่ (25 กิโลเมตรต่อชั่วโมงในเขตเมือง และ 60 กิโลเมตรต่อชั่วโมงเมื่อระยะทางเกิน 15 กิโลเมตร) มิได้ใช้ข้อมูลเส้นทางหรือสภาพจราจรจริง เวลาในตารางจึงอาจคลาดเคลื่อนจากการเดินทางจริง
\item การสร้างคำอธิบายสถานที่ด้วย AI ทำได้ทีละสถานที่จากฟอร์มของผู้ดูแลระบบ หรือระหว่างการนำเข้าตามพื้นที่ได้ครั้งละไม่เกิน 20 แห่ง และเฉพาะสถานที่ที่เพิ่งนำเข้า หากต้องการสร้างคำอธิบายใหม่ให้สถานที่ที่มีอยู่แล้วทั้งหมดในคราวเดียวยังต้องรันสคริปต์แยกนอกระบบ
\end{enumerate}

\subsection{ข้อจำกัดด้านฮาร์ดแวร์และซอฟต์แวร์}

ระบบพึ่งพาการประมวลผลผ่านโมเดลภาษาขนาดใหญ่ตระกูล Gemini ที่เรียกผ่าน Vertex AI ของ Google Cloud ซึ่งกำหนดโควต้าอัตราคำขอต่อโครงการและคิดค่าใช้จ่ายตามปริมาณ Token ทำให้อาจเกิดการหยุดชะงักของบริการเมื่อคำขอเกินโควต้าหากมีผู้ใช้งานพร้อมกันจำนวนมาก และค่าใช้จ่ายเพิ่มขึ้นตามปริมาณการใช้งาน โดยช่วงวันที่ 10--15 ตุลาคม 2569 มีค่าใช้ Vertex AI รวม 135.24 บาท ซึ่งชำระด้วยเครดิตทดลองใช้ฟรีของ Google Cloud (ตารางที่ \ref{tab:cost}) เมื่อเครดิตหมดลงค่าใช้จ่ายส่วนนี้จะเป็นภาระของผู้ดูแลระบบโดยตรง นอกจากนี้ เกณฑ์ของกลไก LLM-as-a-Judge ชุดปัจจุบัน (ผ่านเมื่อไม่มีปัญหาระดับ major) ปรับปรุงจากการตรวจแผนจริงเพียง 10 แผน และการแบ่งปัญหาเป็นระดับ major หรือ minor อาศัยวิจารณญาณของ LLM ซึ่งยังไม่ได้เทียบกับการตัดสินของมนุษย์ และยังไม่ได้วัดอัตราผ่านของเกณฑ์ชุดนี้กับ gemini-3.8-flash อย่างเป็นระบบ

% -----------------------------------------------------------
\section{ปัญหาอุปสรรค และแนวทางแก้ไข}
% -----------------------------------------------------------

ในระหว่างการพัฒนาและทดสอบระบบ คณะผู้จัดทำพบปัญหาทางเทคนิคหลายประการ โดยเฉพาะในส่วนของการควบคุมการใช้งานโควต้าของ LLM และความเร็วในการตอบสนองของแชตบอตและระบบวางแผนการเดินทาง ซึ่งได้มีการวิเคราะห์สาเหตุและดำเนินการแก้ไขบางส่วนไปแล้ว ดังแสดงในตารางที่ \ref{tab:problems}

{\small
\begin{longtable}{|>{\raggedright\arraybackslash}p{2.9cm}|>{\raggedright\arraybackslash}p{4.9cm}|>{\raggedright\arraybackslash}p{5.5cm}|}
\caption{ตารางแสดงการสรุปปัญหา อุปสรรค และแนวทางแก้ไข}\label{tab:problems}\\
\hline
\textbf{ปัญหาและอุปสรรค} & \textbf{สาเหตุหลัก} & \textbf{แนวทางแก้ไข} \\
\hline
\endfirsthead
\caption[]{ตารางแสดงการสรุปปัญหา อุปสรรค และแนวทางแก้ไข~(\textnormal{ต่อ})}\\
\hline
\textbf{ปัญหาและอุปสรรค} & \textbf{สาเหตุหลัก} & \textbf{แนวทางแก้ไข} \\
\hline
\endhead
โควต้าของ LLM อาจหมดลงได้หากมีผู้ใช้งานพร้อมกันจำนวนมาก & ระบบพึ่งพาการประมวลผลผ่าน Gemini บน Vertex AI ซึ่งมีการกำหนดเกณฑ์การใช้งาน (Rate Limits) และคิดค่าใช้จ่ายตามปริมาณ Token & 1. รวมการเรียก LLM ของทุกงานไว้ที่โมดูลกลางเดียว และกำหนดผู้ให้บริการกับโมเดลด้วยตัวแปร LLM\_PROVIDER และ MODEL\_NAME ชุดเดียว จึงสลับไปใช้ผู้ให้บริการอื่นที่รองรับไว้ (KKU AI Gateway หรือ Claude) ได้ทั้งระบบโดยไม่แก้โค้ดเมื่อโควต้าของ Vertex AI มีปัญหา แม้ปัจจุบันทุกงานจะใช้ gemini-3.8-flash \newline 2. จำกัดจำนวนรอบการสร้างแผนซ้ำของ LLM-as-a-Judge ไว้ไม่เกินค่าที่กำหนด (MAX\_JUDGE\_ATTEMPTS) เพื่อไม่ให้คำขอเดียวใช้โควต้าเกินความจำเป็น \newline 3. ตั้งค่าให้ขั้นตอนตรวจสอบคุณภาพ (Judge) ทำงานแบบ Fail-Open หากเรียก API ไม่สำเร็จ เพื่อไม่ให้ปัญหาโควต้าที่ขั้นตอนตรวจสอบไปปิดกั้นผลลัพธ์หลักที่สร้างเสร็จแล้ว \newline 4. จำกัดจำนวนคำขอต่อผู้ใช้แต่ละราย (Rate Limiting) ที่ AI Service โดยค่าเริ่มต้นคือแชตบอต 20 ครั้งต่อนาทีและ 200 ครั้งต่อชั่วโมง ส่วนการวางแผนทริป 5 ครั้งต่อ 10 นาทีและ 20 ครั้งต่อชั่วโมง \newline 5. บันทึกปริมาณ Token และค่าใช้จ่ายโดยประมาณของแต่ละคำขอ เพื่อติดตามการใช้โควต้า \\
\hline
แชตบอตตอบช้า โดยเฉพาะคำถามต่อเนื่องตั้งแต่รอบที่ 2 และได้คะแนนความรวดเร็วในการตอบต่ำที่สุดในแบบสอบถาม (3.87) & ตั้งแต่คำถามรอบที่ 2 ระบบเรียก LLM สองครั้งต่อเนื่องกัน คือเขียนคำถามใหม่ให้เข้าใจได้ด้วยตัวเอง (Query Rewrite) แล้วจึงสร้างคำตอบ ประกอบกับ gemini-3.8-flash ใช้การคิดก่อนตอบ (Reasoning) เป็นค่าเริ่มต้นแม้กับงานสั้น ๆ ทำให้เวลารอเพิ่มขึ้นประมาณ 3--15 วินาทีอย่างไม่แน่นอน & 1. ตั้งระดับการคิด (reasoning\_effort) เป็น low เฉพาะขั้นเขียนคำถามใหม่ ลดเวลาขั้นนี้จากประมาณ 8.5 วินาทีเหลือประมาณ 3 วินาที (ขั้นสร้างคำตอบยังใช้ระดับเดิม เพราะเมื่อลดลงชุดคำถาม Golden Set ไม่ผ่าน 10 จาก 33 ข้อ) \newline 2. ส่งคำตอบแบบทยอยแสดงผล (Streaming) ให้ผู้ใช้เห็นข้อความก่อนคำตอบเสร็จ \newline 3. กำหนดเวลารอ LLM ขั้นเขียนคำถามใหม่ 8 วินาที (เกินแล้วใช้คำถามเดิม) และขั้นสร้างคำตอบ 25 วินาที ลองซ้ำได้ 1 ครั้ง \\
\hline
การสร้างแผนการเดินทางใช้เวลานาน ผู้ใช้ให้คะแนนความรวดเร็วในการสร้างแผนเที่ยว 3.90 และคะแนนข้อนี้กระจายมากที่สุด (S.D. 1.25) & การสร้างแผนต้องเรียก LLM หลายขั้นต่อกัน ได้แก่ การคัดเลือกสถานที่ ซึ่งในการทดลอง Prompt ด้วยโมเดลจริงใช้เวลาเฉลี่ยประมาณ 27.5 วินาทีต่อครั้งสำหรับ Prompt ที่ใช้งานจริง ตามด้วยการตรวจคุณภาพของ Judge และหากแผนไม่ผ่านต้องคัดเลือกใหม่อีกรอบ อีกทั้งไม่สามารถทยอยแสดงผลได้ เพราะตัวจัดตารางต้องได้ผลลัพธ์ JSON ครบทั้งชุดก่อนจึงจัดลำดับและคำนวณเวลาได้ & 1. จำกัดวงรอบสร้างแผนและตรวจคุณภาพไว้ไม่เกิน 2 รอบ (MAX\_JUDGE\_ATTEMPTS) \newline 2. กำหนดเวลารอ LLM ไว้ 120 วินาทีสำหรับการคัดเลือกสถานที่ และ 60 วินาทีสำหรับ Judge โดยปิดการลองซ้ำอัตโนมัติของไลบรารี เพื่อไม่ให้คำขอที่ค้างทำให้ผู้ใช้รอที่หน้ากำลังโหลดเป็นเวลานาน \newline 3. ให้ Judge ทำงานแบบ Fail-Open หากเรียก API ไม่สำเร็จ ระบบส่งแผนที่สร้างเสร็จแล้วกลับไปโดยไม่รอการตรวจ \newline 4. ให้ LLM คืนเฉพาะรหัสสถานที่ของแต่ละวัน ส่วนการจัดลำดับและคำนวณเวลาทำด้วยโค้ดเชิงกำหนดแน่นอน \newline 5. ยังไม่ได้วัดเวลาสร้างแผนทั้งกระบวนการในการใช้งานจริง (ดูหัวข้อข้อเสนอแนะ) \\
\hline
\end{longtable}
\par}

% -----------------------------------------------------------
\section{ข้อเสนอแนะ ในการพัฒนาต่อไป}
% -----------------------------------------------------------

\subsection{การปรับปรุงฟังก์ชันการทำงาน}

เพื่อให้ระบบมีความสมบูรณ์และพร้อมใช้งานจริงยิ่งขึ้น ควรเสริมมาตรการความปลอดภัยของเว็บแอปพลิเคชันที่ยังเปิดอยู่ ได้แก่ การกำหนด Content Security Policy การตรวจชนิดไฟล์อัปโหลดจากเนื้อไฟล์ และการล็อกบัญชีหลังเข้าสู่ระบบผิดซ้ำ ควรจัดให้ฐานข้อมูลที่ใช้งานจริงมีประวัติ migration เพื่อปรับปรุงโครงสร้างได้อัตโนมัติ และเก็บความคืบหน้าของชุดงานนำเข้าไว้ในฐานข้อมูลแทนหน่วยความจำ ควรนำกรณีทดสอบเชิงหน้าที่ของส่วนเว็บเข้าไว้ในชุดทดสอบอัตโนมัติของที่เก็บโค้ด และรันชุดทดสอบของ Backend ใน GitHub Actions ทดสอบฟรอนต์เอนด์ และทดสอบภาระงานเมื่อมีผู้ใช้พร้อมกันจำนวนมาก ควรพัฒนาช่องทางให้ผู้ใช้ปรับแผนการเดินทางหลังสร้างเสร็จ โดยส่งความเห็นของผู้ใช้กลับไปให้ระบบวางแผนประมวลผล นอกจากนี้ ควรศึกษาแนวทางการเข้าถึงข้อมูลความหนาแน่นของผู้คนจากแหล่งข้อมูลเปิดอื่นที่ยังไม่เคยทดลอง เพื่อเพิ่มประสิทธิภาพในการแนะนำแผนการเดินทางให้สอดคล้องกับสถานการณ์จริงยิ่งขึ้นและเมื่อฟังก์ชันหลักครบถ้วนสมบูรณ์ ควรทดสอบยอมรับจากผู้ใช้แบบกำหนดภารกิจ (UAT) และเก็บแบบสอบถามความพึงพอใจจากกลุ่มผู้ใช้ที่กว้างและหลากหลายกว่าเดิม เพื่อประเมินประสิทธิภาพและความพึงพอใจของผู้ใช้งาน รวมถึงเทียบการแบ่งระดับปัญหาของกลไก LLM-as-a-Judge กับการตัดสินของมนุษย์ ทดลอง Prompt ซ้ำกับตัวสร้างแผนฉบับปัจจุบัน ทดลองระบบวางแผนการเดินทางแบบรวมวงรอบสร้างแผนใหม่ของ Judge และประเมินแชตบอตตามตัวชี้วัด Answer Relevance, Faithfulness และ Context Recall/Precision ที่วางแผนไว้ ตลอดจนเปรียบเทียบโมเดลภาษาขนาดใหญ่หลายรุ่นในด้านความแม่นยำ ความเร็วในการตอบสนอง และต้นทุนการใช้งาน เพื่อเลือกโมเดลที่เหมาะสมกับแต่ละฟังก์ชัน (แชตบอต, Trip Planner, Judge) แทนการใช้ gemini-3.8-flash รุ่นเดียวกับทุกฟังก์ชันดังในปัจจุบัน ซึ่งหากผลการเปรียบเทียบชี้ว่าควรใช้ต่างรุ่นกัน ต้องเพิ่มการตั้งค่าโมเดลแยกตามงานกลับเข้าไปในโมดูลเรียก LLM กลาง

\subsection{แนวทางการวิจัยต่อยอด}

นอกจากการปรับปรุงฟังก์ชันข้างต้น ผลของโครงงานนี้ชี้ให้เห็นประเด็นที่ควรศึกษาต่อยอดในเชิงวิจัย ได้แก่ การตรวจสอบความน่าเชื่อถือของการประเมินด้วย LLM โดยเทียบคะแนนของ LLM-as-a-Judge และผู้ประเมิน LLM ในการทดลองกับคะแนนจากผู้ประเมินที่เป็นมนุษย์ ซึ่งการทดลองในบทที่ 4 พบว่าผู้ประเมิน LLM ทั้งสามตัวให้คะแนนรายรันไม่สอดคล้องกัน การนำข้อมูลเส้นทางและสภาพการจราจรจริงมาใช้ในตัวจัดตารางแทนการประมาณระยะทางแบบเส้นตรง เพื่อให้ตารางเวลาในแผนเป็นไปได้จริงมากขึ้น ซึ่งเป็นข้อที่ผู้ใช้ให้คะแนนต่ำในแบบสอบถาม และการวัดผลของระบบต่อการวางแผนท่องเที่ยวจริง เช่น เวลาที่ผู้ใช้ใช้วางแผนเมื่อเทียบกับวิธีเดิม เนื่องจากผลในโครงงานนี้เป็นความเห็นของผู้ใช้ มิใช่ค่าที่วัดจากการใช้งานจริง
